\documentclass[11pt]{article}
\usepackage[margin=1in]{geometry}
\usepackage{amsmath,amssymb,amsthm,mathtools}
\usepackage{graphicx}
\usepackage{microtype}
\usepackage[hidelinks]{hyperref}

\newtheorem{theorem}{Theorem}
\newtheorem{remark}[theorem]{Remark}
\newcommand{\F}{\mathbb F}
\newcommand{\ip}[2]{\langle #1,#2\rangle}
\newcommand{\one}{\mathbf 1}

\title{Rigidity of Optimal Riesz Bounds for Induced Outer Products}
\author{Candidate full solution packet\\\small arXiv:1410.7755, question after Theorem 4.6}
\date{11 August 2026}

\begin{document}
\maketitle

\noindent\textbf{Status.} Candidate full negative answer, pending expert
review: no frames other than equiangular tight frames attain the simultaneous
optimal pair.  More strongly, equality in the sharp lower bound alone forces
an equiangular tight frame.

\section{Source question}

Casazza, Pinkham, and Tuomanen [1] prove that the induced outer products
of a unit-norm frame in $\F^N$ have upper Riesz bound at least $M/N$ and lower
Riesz bound at most $M(N-1)/(N(M-1))$.  Equiangular tight frames attain both.
They then ask whether any other frames can attain these optimal bounds:

\begin{center}
\includegraphics[width=0.98\textwidth]{figures/source_question_crop.png}
\end{center}

Here $\F$ is $\mathbb R$ or $\mathbb C$, and $1<N<M$.  For unit vectors
$\Phi=(\phi_i)_{i=1}^M\subset\F^N$, let
\[
  H=\bigl[|\ip{\phi_i}{\phi_j}|^2\bigr]_{i,j=1}^M
\]
be the Gram matrix of the rank-one operators $\phi_i\phi_i^*$.  Write its
eigenvalues as $\lambda_1\geq\cdots\geq\lambda_M$.  When the outer products
are a Riesz sequence, their optimal bounds are
\[
  B=\lambda_1,\qquad A=\lambda_M.
\]

\section{Proof intuition}

The source's lower-bound proof contains the whole rigidity mechanism.  If the
smallest eigenvalue reaches its upper limit, both the estimate on the largest
eigenvalue and the averaging estimate on the other eigenvalues must be exact.
Thus the outer-product Gram matrix has one eigenvalue of multiplicity one and
one of multiplicity $M-1$.  Its unit diagonal then fixes every coordinate of
the exceptional eigenvector, while entrywise nonnegativity fixes their phases.
This makes all off-diagonal entries equal.  Finally, equality in the frame
potential inequality forces tightness.

\section{Full characterization}

Set
\[
  \beta:=\frac MN,
  \qquad
  \alpha:=\frac{M(N-1)}{N(M-1)},
  \qquad
  c:=\frac{M-N}{N(M-1)}.
\]

\begin{theorem}\label{thm:main}
Let $\Phi=(\phi_i)_{i=1}^M$ be a unit-norm frame for $\F^N$, where
$1<N<M$, and suppose its induced outer products form a Riesz sequence with
optimal bounds $A\leq B$.  The following are equivalent:
\begin{enumerate}
  \item $A=\alpha$;
  \item $(A,B)=(\alpha,\beta)$;
  \item $B-A=\beta-\alpha$;
  \item $\Phi$ is an equiangular tight frame.
\end{enumerate}
In this case $|\ip{\phi_i}{\phi_j}|^2=c$ whenever $i\ne j$.
Consequently, there are no other frames attaining the optimal pair asked
about in [1].
\end{theorem}

\begin{proof}
The matrix $H$ is positive semidefinite and has diagonal entries one, so
$\operatorname{tr}H=M$.  Its entry sum is the frame potential.  The standard
frame-potential inequality and the Rayleigh quotient give
\[
 \frac{M^2}{N}
 \leq \one^*H\one
 \leq M\lambda_1,
 \qquad\text{hence}\qquad
 \lambda_1\geq\beta.
\]
Moreover,
\[
 (M-1)\lambda_M
 \leq\sum_{i=2}^M\lambda_i
 =M-\lambda_1
 \leq M-\beta=(M-1)\alpha,
\]
so $A=\lambda_M\leq\alpha$.

Assume first that $A=\alpha$.  The last displayed chain has equal endpoints,
so every inequality in it is equality.  Therefore
\[
 \lambda_1=\beta,
 \qquad
 \lambda_2=\cdots=\lambda_M=\alpha.
\]
Since $\beta>\alpha$, the spectral theorem supplies a unit vector
$u\in\F^M$ such that
\[
 H=\alpha I_M+(\beta-\alpha)uu^*.
\]
Taking diagonal entries and simplifying gives, for every $i$,
\[
 |u_i|^2=\frac{1-\alpha}{\beta-\alpha}=\frac1M.
\]
For $i\ne j$, the number
\[
 H_{ij}=(\beta-\alpha)u_i\overline{u_j}
       =|\ip{\phi_i}{\phi_j}|^2
\]
is real and nonnegative, while its modulus is
\[
 |H_{ij}|=\frac{\beta-\alpha}{M}
          =\frac{M-N}{N(M-1)}=c.
\]
Thus $H_{ij}=c$ for all $i\ne j$, proving equiangularity.

It remains to prove tightness without assuming it.  If $S$ is the frame
operator, then
\[
 \operatorname{tr}S=M,
 \qquad
 \operatorname{tr}(S^2)
   =\sum_{i,j=1}^M|\ip{\phi_i}{\phi_j}|^2
   =M+M(M-1)c=\frac{M^2}{N}.
\]
For the $N$ nonnegative eigenvalues of $S$, Cauchy--Schwarz says
$\operatorname{tr}(S^2)\geq(\operatorname{tr}S)^2/N$, with equality only
when they are all $M/N$.  Hence $S=(M/N)I_N$, and $\Phi$ is tight.  This
proves $(1)\Rightarrow(4)$ and also yields $(2)$.

Conversely, a unit-norm equiangular tight frame meets the Welch equality, so
its squared common angle is $c$.  Consequently
\[
 H=(1-c)I_M+c\one\one^*,
\]
whose eigenvalues are $1+(M-1)c=\beta$ once and
$1-c=\alpha$ with multiplicity $M-1$.  Thus $(4)\Rightarrow(2)$, and
$(2)\Rightarrow(1)$ is immediate.

Finally, the two universal inequalities give
\[
 B-A-(\beta-\alpha)=(B-\beta)+(\alpha-A)\geq0.
\]
Equality holds precisely when both nonnegative terms vanish, which is
equivalent to $(2)$.  This proves the equivalence with $(3)$.
\end{proof}

\section{Scope and edge cases}

The theorem distinguishes the two one-sided questions.  Equality in the
\emph{upper} bound alone characterizes all unit-norm tight frames, as already
proved in [1]; it does not by itself force equiangularity.  In contrast,
equality in the \emph{lower} bound already forces the complete equiangular
tight structure.  Since $\alpha>0$, attaining it automatically makes $H$
positive definite, so the Riesz assumption may equivalently be omitted from
the hypothesis and recovered from the equality statement.

The restriction $M>N>1$ is the redundant, nondegenerate regime of the source
question.  Existence of equality still depends on existence of an
equiangular tight frame with parameters $(N,M)$; when no such frame exists,
the universal numerical bound is not attained.

\section{Verification and provenance}

The proof uses only the Gram identity from [1], the trace and Rayleigh
quotient, the frame-potential identity, and equality in Cauchy--Schwarz.  A
constant-check script bundled with this packet verifies all rational
identities over a broad finite parameter range.  Cheap run indexes were
searched for the exact title, authors, arXiv id, and the core phrases
``outer-product Riesz bounds'' and ``equiangular outer products.''  No later
answer was found; the sole later parsed arXiv citation located was unrelated
to this equality question.  Novelty therefore remains subject to specialist
review.

\medskip
\noindent\textbf{Reference.} [1]
P. G. Casazza, E. Pinkham, and B. Tuomanen,
\emph{Riesz outer product Hilbert space frames: quantitative bounds,
topological properties, and full geometric characterization}, J. Math. Anal.
Appl. 441 (2016), 475--498; arXiv:1410.7755.

\end{document}
